% =============================================================================
% ANNEX F : TEAM RETROSPECTIVE REPORT (B4), REDACTED
% =============================================================================

\chapter{First team retrospective report (B4)}
\label{annex:retro-b4}

\begin{tcolorbox}[colback=blue!5, colframe=blue!40!black, boxrule=0.5pt, arc=3pt,
  left=8pt, right=8pt, top=6pt, bottom=6pt]
  \small\itshape
  \textbf{Avertissement.} Cette annexe est une reproduction \textbf{non fidele} d'un compte-rendu de retrospective utilise en contexte professionnel. Elle est volontairement adaptee, tronquee et caviardee a des fins de \textbf{demonstration} et de \textbf{preuve} dans le portfolio. Les informations sensibles (identites, dates exactes, identifiants, references internes et metadonnees d'outillage) ont ete masquees.
\end{tcolorbox}

\bigskip

\begin{tcolorbox}[colback=gray!8, colframe=gray!50, boxrule=0.5pt, arc=3pt,
  left=8pt, right=8pt, top=6pt, bottom=6pt]
\begin{tabularx}{\linewidth}{@{}p{4.3cm}X@{}}
\textbf{Team} & TWIICE software team \\
\textbf{Related sprint} & Sprint \caviarder[1.8cm] (2 weeks) \\
\textbf{Retrospective date} & \caviarder[2.2cm] \\
\textbf{Facilitator} & \caviarder[2.6cm] \\
\textbf{Participants} & \caviarder[6cm] \\
\textbf{Context} & First complete retrospective after deployment of an IEC~62304-adapted agile framework \\
\end{tabularx}
\end{tcolorbox}

\section*{1. Retrospective objective}

Evaluate the sprint across three axes: what should continue, what should stop, and what should start in the next sprint. The operational goal is to improve agile delivery pace while strengthening certification-grade documentation quality.

\section*{2. Method applied}

The session follows a Start/Stop/Continue format, with timeboxing, collective priority voting, and final action validation. Each retained action is assigned an owner, a deadline, and a verification criterion.

\section*{3. Start / Stop / Continue results}

\begin{tabularx}{\textwidth}{|p{2.2cm}|X|}
\hline
\textbf{Continue} &
\begin{itemize}
  \item Maintain the 15-minute daily ritual focused on technical blockers and compliance risks.
  \item Keep the value $\times$ certification-risk prioritization model in the backlog.
  \item Continue cross-review between development and quality before sprint closure.
\end{itemize} \\
\hline
\textbf{Stop} &
\begin{itemize}
  \item Stop late end-of-sprint merges without complete documentation checks.
  \item Stop ambiguous tickets without explicit acceptance criteria.
  \item Stop ad hoc anomaly handling outside the prioritized queue.
\end{itemize} \\
\hline
\textbf{Start} &
\begin{itemize}
  \item Start a single DoD checklist (code, tests, IEC~62304 docs, quality review).
  \item Start a standardized "residual risks" section in each decision note.
  \item Start weekly monitoring of traceability gaps (requirements $\rightarrow$ tests $\rightarrow$ evidence).
\end{itemize} \\
\hline
\end{tabularx}

\section*{4. Retained actions and execution evidence}

\begin{tabularx}{\textwidth}{|p{4.6cm}|p{2.7cm}|p{2.4cm}|X|}
\hline
\textbf{Action} & \textbf{Owner} & \textbf{Deadline} & \textbf{Verification criterion} \\
\hline
Common DoD checklist activated on all tickets in the next sprint & \caviarder[2.1cm] & \caviarder[1.8cm] & 100\% of "Done" tickets include code + tests + documented traceability \\
\hline
"Residual risks" template added to decision notes & \caviarder[2.1cm] & \caviarder[1.8cm] & Each new decision note includes risks + mitigation section \\
\hline
Weekly traceability review institutionalized & \caviarder[2.1cm] & \caviarder[1.8cm] & Gap dashboard maintained; critical gaps closed before gate review \\
\hline
\end{tabularx}

\section*{5. How this evidence demonstrates training impact}

This retrospective evidences a change in maturity level.

\begin{itemize}
  \item \textbf{Before training}: verbal feedback loops, weak traceability, and limited action follow-through.
  \item \textbf{After training}: explicit improvement loop (observation $\rightarrow$ action $\rightarrow$ verification) with auditable artifacts.
\end{itemize}

The observable benefit is twofold: team dynamics become more stable (fewer repeated debates), and compliance becomes integrated into daily steering rather than appended at end of cycle.

\section*{6. Closure status}

\begin{tabularx}{\textwidth}{|p{4.2cm}|X|}
\hline
\textbf{Team validation} & Retrospective approved by the project team (redacted trace) \\
\hline
\textbf{Documentation trace} & Record linked to sprint artifacts \caviarder[2.0cm] and to gate preparation \caviarder[1.6cm] \\
\hline
\textbf{Link to B4 competency} & Evidence of genuine iterative adaptation in a regulated environment \\
\hline
\end{tabularx}


